% jbr-header.tex (jbrquarto 1.0.0)
% Preamble for Journal of Business Research (JBR) manuscripts. pdfLaTeX fonts: Times for the text,
% tables, captions, and notes; Helvetica for any sans serif text.

% ---- Fonts ----
\usepackage{mathptmx}
\usepackage[scaled=0.92]{helvet}

% ---- Quarto 1.4 (pandoc before 3.2) ----
% Those versions wrap \paragraph and \subparagraph in a way that titlesec
% cannot read ("Argument of \paragraph has an extra }"). The original
% commands are put back before titlesec is loaded; the heading formats below
% define them anew.
\makeatletter
\@ifundefined{xxxParagraphStar}{%
  \@ifundefined{oldparagraph}{}{\let\paragraph\oldparagraph}%
  \@ifundefined{oldsubparagraph}{}{\let\subparagraph\oldsubparagraph}%
}{}
\makeatother

% ---- Packages ----
\usepackage{float}
\usepackage{setspace}
\usepackage{ragged2e}
\usepackage{titlesec}
\usepackage{etoolbox}
\usepackage{longtable}
\usepackage{booktabs}
\usepackage[bottom]{footmisc}
\usepackage{indentfirst}
\usepackage{fancyhdr}

% ---- Characters typed directly ----
% pdfLaTeX has no glyphs of its own for Greek letters and for many signs that
% are typed as characters (a Greek alpha, a minus sign, "less than or equal").
% In the text, Greek letters are taken upright from the Symbol font, as they
% appear in a Word file; in formulas they are the usual math letters. The
% other signs are set in math mode. XeLaTeX and LuaLaTeX read all of these
% characters without help, so the definitions apply to pdfLaTeX only.
\usepackage{iftex}
\usepackage{pifont}
\ifPDFTeX
\providecommand{\texorpdfstring}[2]{#1}
\newcommand{\jbrgreek}[3]{\texorpdfstring{\TextOrMath{\Pisymbol{psy}{#1}}{#2}}{#3}}
\newcommand{\jbrsign}[1]{\texorpdfstring{\ensuremath{#1}}{}}
\DeclareUnicodeCharacter{03B1}{\jbrgreek{97}{\alpha}{alpha}}
\DeclareUnicodeCharacter{03B2}{\jbrgreek{98}{\beta}{beta}}
\DeclareUnicodeCharacter{03B3}{\jbrgreek{103}{\gamma}{gamma}}
\DeclareUnicodeCharacter{03B4}{\jbrgreek{100}{\delta}{delta}}
\DeclareUnicodeCharacter{03B5}{\jbrgreek{101}{\varepsilon}{epsilon}}
\DeclareUnicodeCharacter{03F5}{\jbrgreek{101}{\epsilon}{epsilon}}
\DeclareUnicodeCharacter{03B6}{\jbrgreek{122}{\zeta}{zeta}}
\DeclareUnicodeCharacter{03B7}{\jbrgreek{104}{\eta}{eta}}
\DeclareUnicodeCharacter{03B8}{\jbrgreek{113}{\theta}{theta}}
\DeclareUnicodeCharacter{03D1}{\jbrgreek{74}{\vartheta}{theta}}
\DeclareUnicodeCharacter{03B9}{\jbrgreek{105}{\iota}{iota}}
\DeclareUnicodeCharacter{03BA}{\jbrgreek{107}{\kappa}{kappa}}
\DeclareUnicodeCharacter{03BB}{\jbrgreek{108}{\lambda}{lambda}}
\DeclareUnicodeCharacter{03BC}{\jbrgreek{109}{\mu}{mu}}
\DeclareUnicodeCharacter{03BD}{\jbrgreek{110}{\nu}{nu}}
\DeclareUnicodeCharacter{03BE}{\jbrgreek{120}{\xi}{xi}}
\DeclareUnicodeCharacter{03BF}{\jbrgreek{111}{o}{o}}
\DeclareUnicodeCharacter{03C0}{\jbrgreek{112}{\pi}{pi}}
\DeclareUnicodeCharacter{03D6}{\jbrgreek{118}{\varpi}{pi}}
\DeclareUnicodeCharacter{03C1}{\jbrgreek{114}{\rho}{rho}}
\DeclareUnicodeCharacter{03C2}{\jbrgreek{86}{\varsigma}{sigma}}
\DeclareUnicodeCharacter{03C3}{\jbrgreek{115}{\sigma}{sigma}}
\DeclareUnicodeCharacter{03C4}{\jbrgreek{116}{\tau}{tau}}
\DeclareUnicodeCharacter{03C5}{\jbrgreek{117}{\upsilon}{upsilon}}
\DeclareUnicodeCharacter{03C6}{\jbrgreek{106}{\varphi}{phi}}
\DeclareUnicodeCharacter{03D5}{\jbrgreek{102}{\phi}{phi}}
\DeclareUnicodeCharacter{03C7}{\jbrgreek{99}{\chi}{chi}}
\DeclareUnicodeCharacter{03C8}{\jbrgreek{121}{\psi}{psi}}
\DeclareUnicodeCharacter{03C9}{\jbrgreek{119}{\omega}{omega}}
\DeclareUnicodeCharacter{0391}{\jbrgreek{65}{\mathrm{A}}{A}}
\DeclareUnicodeCharacter{0392}{\jbrgreek{66}{\mathrm{B}}{B}}
\DeclareUnicodeCharacter{0393}{\jbrgreek{71}{\Gamma}{Gamma}}
\DeclareUnicodeCharacter{0394}{\jbrgreek{68}{\Delta}{Delta}}
\DeclareUnicodeCharacter{2206}{\jbrgreek{68}{\Delta}{Delta}}
\DeclareUnicodeCharacter{0395}{\jbrgreek{69}{\mathrm{E}}{E}}
\DeclareUnicodeCharacter{0396}{\jbrgreek{90}{\mathrm{Z}}{Z}}
\DeclareUnicodeCharacter{0397}{\jbrgreek{72}{\mathrm{H}}{H}}
\DeclareUnicodeCharacter{0398}{\jbrgreek{81}{\Theta}{Theta}}
\DeclareUnicodeCharacter{0399}{\jbrgreek{73}{\mathrm{I}}{I}}
\DeclareUnicodeCharacter{039A}{\jbrgreek{75}{\mathrm{K}}{K}}
\DeclareUnicodeCharacter{039B}{\jbrgreek{76}{\Lambda}{Lambda}}
\DeclareUnicodeCharacter{039C}{\jbrgreek{77}{\mathrm{M}}{M}}
\DeclareUnicodeCharacter{039D}{\jbrgreek{78}{\mathrm{N}}{N}}
\DeclareUnicodeCharacter{039E}{\jbrgreek{88}{\Xi}{Xi}}
\DeclareUnicodeCharacter{039F}{\jbrgreek{79}{\mathrm{O}}{O}}
\DeclareUnicodeCharacter{03A0}{\jbrgreek{80}{\Pi}{Pi}}
\DeclareUnicodeCharacter{03A1}{\jbrgreek{82}{\mathrm{P}}{P}}
\DeclareUnicodeCharacter{03A3}{\jbrgreek{83}{\Sigma}{Sigma}}
\DeclareUnicodeCharacter{03A4}{\jbrgreek{84}{\mathrm{T}}{T}}
\DeclareUnicodeCharacter{03A5}{\jbrgreek{85}{\Upsilon}{Upsilon}}
\DeclareUnicodeCharacter{03A6}{\jbrgreek{70}{\Phi}{Phi}}
\DeclareUnicodeCharacter{03A7}{\jbrgreek{67}{\mathrm{X}}{X}}
\DeclareUnicodeCharacter{03A8}{\jbrgreek{89}{\Psi}{Psi}}
\DeclareUnicodeCharacter{03A9}{\jbrgreek{87}{\Omega}{Omega}}
\DeclareUnicodeCharacter{2212}{\texorpdfstring{\ensuremath{-}}{-}}
\DeclareUnicodeCharacter{2264}{\jbrsign{\leq}}
\DeclareUnicodeCharacter{2265}{\jbrsign{\geq}}
\DeclareUnicodeCharacter{2260}{\jbrsign{\neq}}
\DeclareUnicodeCharacter{2248}{\jbrsign{\approx}}
\DeclareUnicodeCharacter{2261}{\jbrsign{\equiv}}
\DeclareUnicodeCharacter{223C}{\jbrsign{\sim}}
\DeclareUnicodeCharacter{2245}{\jbrsign{\cong}}
\DeclareUnicodeCharacter{221D}{\jbrsign{\propto}}
\DeclareUnicodeCharacter{2213}{\jbrsign{\mp}}
\DeclareUnicodeCharacter{2217}{\jbrsign{\ast}}
\DeclareUnicodeCharacter{22C5}{\jbrsign{\cdot}}
\DeclareUnicodeCharacter{2219}{\jbrsign{\bullet}}
\DeclareUnicodeCharacter{221A}{\jbrsign{\surd}}
\DeclareUnicodeCharacter{221E}{\jbrsign{\infty}}
\DeclareUnicodeCharacter{2211}{\jbrsign{\sum}}
\DeclareUnicodeCharacter{220F}{\jbrsign{\prod}}
\DeclareUnicodeCharacter{222B}{\jbrsign{\int}}
\DeclareUnicodeCharacter{2202}{\jbrsign{\partial}}
\DeclareUnicodeCharacter{2207}{\jbrsign{\nabla}}
\DeclareUnicodeCharacter{2208}{\jbrsign{\in}}
\DeclareUnicodeCharacter{2209}{\jbrsign{\notin}}
\DeclareUnicodeCharacter{2282}{\jbrsign{\subset}}
\DeclareUnicodeCharacter{2286}{\jbrsign{\subseteq}}
\DeclareUnicodeCharacter{2229}{\jbrsign{\cap}}
\DeclareUnicodeCharacter{222A}{\jbrsign{\cup}}
\DeclareUnicodeCharacter{2205}{\jbrsign{\emptyset}}
\DeclareUnicodeCharacter{2200}{\jbrsign{\forall}}
\DeclareUnicodeCharacter{2203}{\jbrsign{\exists}}
\DeclareUnicodeCharacter{2227}{\jbrsign{\wedge}}
\DeclareUnicodeCharacter{2228}{\jbrsign{\vee}}
\DeclareUnicodeCharacter{2032}{\jbrsign{{}^{\prime}}}
\DeclareUnicodeCharacter{2033}{\jbrsign{{}^{\prime\prime}}}
\DeclareUnicodeCharacter{2194}{\jbrsign{\leftrightarrow}}
\DeclareUnicodeCharacter{21D2}{\jbrsign{\Rightarrow}}
\DeclareUnicodeCharacter{21D0}{\jbrsign{\Leftarrow}}
\DeclareUnicodeCharacter{21D4}{\jbrsign{\Leftrightarrow}}
\DeclareUnicodeCharacter{2197}{\jbrsign{\nearrow}}
\DeclareUnicodeCharacter{2198}{\jbrsign{\searrow}}
\DeclareUnicodeCharacter{25CF}{\jbrsign{\bullet}}
\DeclareUnicodeCharacter{25CB}{\jbrsign{\circ}}
\DeclareUnicodeCharacter{25A0}{\jbrsign{\blacksquare}}
\DeclareUnicodeCharacter{25A1}{\jbrsign{\square}}
\DeclareUnicodeCharacter{25B2}{\jbrsign{\blacktriangle}}
\DeclareUnicodeCharacter{25B3}{\jbrsign{\triangle}}
\DeclareUnicodeCharacter{25BC}{\jbrsign{\blacktriangledown}}
\DeclareUnicodeCharacter{2605}{\jbrsign{\bigstar}}
\DeclareUnicodeCharacter{2606}{\texorpdfstring{\ding{73}}{}}
\DeclareUnicodeCharacter{2023}{\jbrsign{\triangleright}}
\DeclareUnicodeCharacter{2713}{\texorpdfstring{\ding{51}}{}}
\DeclareUnicodeCharacter{2714}{\texorpdfstring{\ding{52}}{}}
\DeclareUnicodeCharacter{2717}{\texorpdfstring{\ding{55}}{}}
\DeclareUnicodeCharacter{2718}{\texorpdfstring{\ding{56}}{}}
\DeclareUnicodeCharacter{2070}{\textsuperscript{0}}
\DeclareUnicodeCharacter{2074}{\textsuperscript{4}}
\DeclareUnicodeCharacter{2075}{\textsuperscript{5}}
\DeclareUnicodeCharacter{2076}{\textsuperscript{6}}
\DeclareUnicodeCharacter{2077}{\textsuperscript{7}}
\DeclareUnicodeCharacter{2078}{\textsuperscript{8}}
\DeclareUnicodeCharacter{2079}{\textsuperscript{9}}
\DeclareUnicodeCharacter{207A}{\textsuperscript{+}}
\DeclareUnicodeCharacter{207B}{\textsuperscript{\ensuremath{-}}}
\DeclareUnicodeCharacter{2080}{\textsubscript{0}}
\DeclareUnicodeCharacter{2081}{\textsubscript{1}}
\DeclareUnicodeCharacter{2082}{\textsubscript{2}}
\DeclareUnicodeCharacter{2083}{\textsubscript{3}}
\DeclareUnicodeCharacter{2084}{\textsubscript{4}}
\DeclareUnicodeCharacter{2085}{\textsubscript{5}}
\DeclareUnicodeCharacter{2086}{\textsubscript{6}}
\DeclareUnicodeCharacter{2087}{\textsubscript{7}}
\DeclareUnicodeCharacter{2088}{\textsubscript{8}}
\DeclareUnicodeCharacter{2089}{\textsubscript{9}}
\DeclareUnicodeCharacter{2153}{\textsuperscript{1}/\textsubscript{3}}
\DeclareUnicodeCharacter{2154}{\textsuperscript{2}/\textsubscript{3}}
\DeclareUnicodeCharacter{2060}{\nobreak}
\DeclareUnicodeCharacter{FB00}{ff}
\DeclareUnicodeCharacter{FB01}{fi}
\DeclareUnicodeCharacter{FB02}{fl}
\DeclareUnicodeCharacter{FB03}{ffi}
\DeclareUnicodeCharacter{FB04}{ffl}
\fi

% ---- Tables: single spaced, in the text font ----
% \setstretch{1} rather than \singlespacing, which adds a blank line
\AtBeginEnvironment{table}{\normalsize\setstretch{1}}
\AtBeginEnvironment{table*}{\normalsize\setstretch{1}}
\AtBeginEnvironment{longtable}{\normalsize\setstretch{1}}
% Space above a table title, and less space between a table and its note
\setlength{\LTpre}{18pt}
\setlength{\LTpost}{6pt}
% A table that fits on a page is not split: \jbrkeep{height} starts a new
% page when less than that height is left on the current page
\newcommand{\jbrkeep}[1]{%
  \par\ifdim\pagetotal>0pt\ifdim\dimexpr\pagegoal-\pagetotal\relax<#1\relax\newpage\fi\fi}

% A table or a figure without a label takes no number. LaTeX counts every
% title, and pandoc 3.2 (in Quarto 1.5) counts every table, with or without
% one; the counter is put back afterward, so that the numbers agree with the
% Word file, and a title is printed without "Table 1" before it.
\newcounter{jbrsaved}
\newcounter{jbrplain}
\newcommand{\jbrplainbegin}[1]{\setcounter{jbrsaved}{\value{#1}}\stepcounter{jbrplain}%
  \begingroup\ifdefined\captionsetup\captionsetup{format=plain,labelformat=empty,labelsep=none}\fi%
  \expandafter\def\csname theH#1\endcsname{plain.\arabic{jbrplain}}}
\newcommand{\jbrplainend}[1]{\endgroup\setcounter{#1}{\value{jbrsaved}}}

% ---- Captions ----
% "Table 1. Title" above a table and "Fig. 1. Caption." below a figure,
% flush left, with the label in bold
\usepackage[labelsep=period,labelfont={bf},textfont={},font={rm,normalsize,singlespacing},
            justification=raggedright,singlelinecheck=false,skip=6pt]{caption}
\captionsetup[table]{position=above}
\captionsetup[figure]{position=below}
% The captions of the panels of a figure: "(a) Text", centered, in the type
% of the text. (Quarto loads the subcaption package when a document has a
% figure with panels.)
\makeatletter
\AtBeginDocument{\@ifpackageloaded{subcaption}{%
  \captionsetup[sub]{font={rm,normalsize,singlespacing},labelfont={},textfont={},
    labelformat=parens,labelsep=space,justification=centering,singlelinecheck=false,skip=4pt}}{}}
\makeatother
% The caption of a picture in a panel without a number: the same, without a letter
\newcommand{\jbrpanelcaption}[1]{{\captionsetup{font={rm,normalsize,singlespacing},labelfont={},
    textfont={},justification=centering,singlelinecheck=false,skip=4pt}\caption*{#1}}}

% ---- No underlining ----
\renewcommand{\underline}[1]{#1}

% ---- Notes under tables and figures ----
% Written in the .qmd file as a ::: {.notes} div; the filter adds the label
\newenvironment{jbrnotes}{%
  \par\nopagebreak[4]\vskip 4pt%
  \begingroup\normalsize\setstretch{1}\raggedright
  \setlength{\parindent}{0pt}\setlength{\parskip}{2pt}%
}{%
  \par\endgroup\vskip 6pt%
}

% ---- Hypotheses: indented, with 12 pt of space above and below ----
% Written in the .qmd file as a ::: {.hypothesis} div; the filter sets the
% label ("H1:") in bold. \addvspace keeps the space single between two
% hypotheses in a row.
\newenvironment{jbrhypothesis}{%
  \par\addvspace{12pt}\begingroup\setlength{\parindent}{0pt}%
  \setlength{\leftskip}{0.5in}%
}{%
  \par\endgroup\addvspace{12pt}%
}

% ---- Centered text, unindented text, and "insert table here" lines ----
\newenvironment{jbrcentered}{\par\begingroup\centering\setlength{\parindent}{0pt}}{\par\endgroup}
\newenvironment{jbrflush}{\par\begingroup\setlength{\parindent}{0pt}}{\par\endgroup}
\newenvironment{jbrcallout}{%
  \par\addvspace{0.5\baselineskip}\begingroup\centering\setlength{\parindent}{0pt}%
}{%
  \par\endgroup\addvspace{0.5\baselineskip}%
}

% ---- References: half-inch hanging indent ----
\ifdefined\cslhangindent
  \setlength{\cslhangindent}{0.5in}
\else
  \newlength{\cslhangindent}
  \setlength{\cslhangindent}{0.5in}
\fi
\newenvironment{jbrreferences}{\par\begingroup}{\par\endgroup}

% ---- Statistics helpers ----
\newcommand{\Mstat}[1]{\ensuremath{M = #1}}
\newcommand{\SDstat}[1]{\ensuremath{\text{SD} = #1}}
\newcommand{\pval}[1]{\ensuremath{p = #1}}
\newcommand{\Fstat}[3]{\ensuremath{F(#1, #2) = #3}}
\newcommand{\tstat}[2]{\ensuremath{t(#1) = #2}}
\newcommand{\chisq}[2]{\ensuremath{\chi^2(#1) = #2}}
\newcommand{\CI}[2]{95\% CI [#1, #2]}

% ---- Front matter (the filter places the text in these environments) ----
% Title and authors
\newenvironment{jbrtitle}{%
  \par\begingroup\centering\bfseries\setlength{\parindent}{0pt}%
}{\par\endgroup}
\newenvironment{jbrauthors}{%
  \par\begingroup\centering\setlength{\parindent}{0pt}%
}{\par\endgroup}
% "Abstract": set like a section heading, without a number
\newenvironment{jbrfronthead}{%
  \par\addvspace{\baselineskip}\begingroup\bfseries\setlength{\parindent}{0pt}%
}{\par\endgroup}
% Abstract, keywords, and the corresponding author: flush left
\newenvironment{jbrabstract}{\par\begingroup\setlength{\parindent}{0pt}}{\par\endgroup}
\newenvironment{jbrkeywords}{\par\begingroup\setlength{\parindent}{0pt}}{\par\endgroup}
\newenvironment{jbrfrontlines}{%
  \par\addvspace{\baselineskip}\begingroup\setlength{\parindent}{0pt}%
}{\par\endgroup}

% ---- MANUSCRIPT SETUP (submission format) ----
% Double spaced (linestretch), 12 pt, 1-inch margins, pages numbered.
% Sections are numbered 1., 1.1., 1.1.1.; every heading stands on its own
% line.
\newcommand{\jbrsetupmanuscript}{%
  \pagestyle{plain}%
  \raggedright%
  \setlength{\parindent}{0.5in}%
  \titleformat{\section}%
    {\normalfont\normalsize\bfseries\raggedright}{\thesection.}{0.5em}{}%
  \titlespacing*{\section}{0pt}{\baselineskip}{0pt}%
  \titleformat{\subsection}%
    {\normalfont\normalsize\itshape\raggedright}{\thesubsection.}{0.5em}{}%
  \titlespacing*{\subsection}{0pt}{0pt}{0pt}%
  \titleformat{\subsubsection}%
    {\normalfont\normalsize\itshape\raggedright}{\thesubsubsection.}{0.5em}{}%
  \titlespacing*{\subsubsection}{0pt}{0pt}{0pt}%
  \titleformat{\paragraph}%
    {\normalfont\normalsize\itshape\raggedright}{\theparagraph.}{0.5em}{}%
  \titlespacing*{\paragraph}{0pt}{0pt}{0pt}%
  \titleformat{\subparagraph}%
    {\normalfont\normalsize\itshape\raggedright}{\thesubparagraph.}{0.5em}{}%
  \titlespacing*{\subparagraph}{0pt}{0pt}{0pt}%
}

% ---- ARTICLE SETUP (publication-style preview; not for submission) ----
\newcommand{\jbrsetuparticle}{%
  \pagestyle{fancy}%
  \fancyhf{}%
  \fancyfoot[C]{\small\thepage}%
  \renewcommand{\headrulewidth}{0pt}%
  \fancypagestyle{plain}{%
    \fancyhf{}%
    \fancyfoot[C]{\small\thepage}%
    \renewcommand{\headrulewidth}{0pt}%
  }%
  \thispagestyle{plain}%
  \setlength{\parindent}{1.5em}%
  \titleformat{\section}%
    {\normalfont\large\bfseries\raggedright}{}{0pt}{}%
  \titlespacing*{\section}{0pt}{18pt}{6pt}%
  \titleformat{\subsection}%
    {\normalfont\normalsize\itshape\raggedright}{}{0pt}{}%
  \titlespacing*{\subsection}{0pt}{14pt}{4pt}%
  \titleformat{\subsubsection}%
    {\normalfont\normalsize\itshape\raggedright}{}{0pt}{}%
  \titlespacing*{\subsubsection}{0pt}{10pt}{3pt}%
  \titleformat{\paragraph}%
    {\normalfont\normalsize\itshape\raggedright}{}{0pt}{}%
  \titlespacing*{\paragraph}{0pt}{10pt}{3pt}%
  \titleformat{\subparagraph}%
    {\normalfont\normalsize\itshape\raggedright}{}{0pt}{}%
  \titlespacing*{\subparagraph}{0pt}{10pt}{3pt}%
  \captionsetup{font=small}%
}
\newenvironment{jbrarttitle}{%
  \par\begingroup\fontsize{16}{19}\selectfont\bfseries\raggedright\setlength{\parindent}{0pt}%
}{\par\endgroup\vskip 10pt}
\newenvironment{jbrartauthors}{%
  \par\begingroup\centering\normalsize\bfseries\setlength{\parindent}{0pt}%
}{\par\endgroup\vskip 6pt}
\newenvironment{jbrartabstract}{%
  \par\noindent\textbf{Abstract}\par\vskip 4pt\begingroup\small\setlength{\parindent}{0pt}%
}{\par\endgroup\vskip 8pt}
\newenvironment{jbrartkeywords}{%
  \par\begingroup\small\setlength{\parindent}{0pt}%
}{\par\endgroup\vskip 8pt}
\newcommand{\jbrartrule}{%
  \vskip 0.8\baselineskip\noindent\rule{\textwidth}{0.4pt}\vskip 0.8\baselineskip}
\makeatletter
\newcommand{\jbrauthornote}[1]{%
  \begingroup\def\@makefnmark{}%
  \footnotetext{\fontsize{8}{10}\selectfont\raggedright #1}\endgroup}
\makeatother
